\documentclass[12pt,notitlepage]{report}
% The shared notes preamble lives outside this repo.
\input{/Users/aryankulkarni/Documents/y3-notes/preamble.tex}

\title{Week 4: binomial regression for injury}
\author{Aryan Kulkarni}
\date{MTH6134 group project}

\begin{document}
\maketitle

\paragraph{Findings.}
Injury risk rises with previous injuries (odds ratio 1.69 per injury) and
with training load, and the effect of load depends on age: there is no evidence of it at
20, and it doubles the odds per point at 30. Position, BMI and club show no
additional effect.

\section{Observed injury rates}

Of the 620 players, 115 (18.5\%) were injured during the season.

\begin{minted}{r}
source("../common.R")
players <- players |>
  mutate(
    age_c = age - 25,
    load_c = training_load - 3,
    age_band = cut(age, c(16, 22, 25, 28, 36),
                   labels = c("17-22", "23-25", "26-28", "29-36")),
    load_band = cut(training_load, c(0, 2, 3, 5), labels = c("1-2", "3", "4-5"))
  )

rates <- players |>
  mutate(across(c(age_band, training_load, prev_injuries), as.character)) |>
  pivot_longer(c(age_band, training_load, prev_injuries),
               names_to = "variable", values_to = "level") |>
  group_by(variable, level) |>
  summarise(n = n(), rate = mean(injured),
            se = sqrt(rate * (1 - rate) / n), .groups = "drop")

fig <- ggplot(rates, aes(x = level, y = rate)) +
  geom_pointrange(aes(ymin = pmax(rate - 1.96 * se, 0),
                      ymax = rate + 1.96 * se),
                  colour = blue, size = 0.3) +
  facet_wrap(~variable, scales = "free_x", labeller = as_labeller(c(
    age_band = "age", training_load = "training load",
    prev_injuries = "previous injuries"
  ))) +
  labs(x = NULL, y = "Proportion injured")
save_fig(fig, "rates", height = 2.5)

out(round(with(players, tapply(injured, list(age = age_band, load = load_band),
                               mean)), 2), "rates-age-load")
\end{minted}

\begin{figure}[H]
  \centering
  \incfig{rates}
  \caption{Proportion injured with 95\% intervals. Risk climbs with age,
    with training load and most steeply with previous injuries.}
\end{figure}

\inputminted{text}{output/rates-age-load.txt}

The two-way table shows more than the margins do. Among players aged 22 or
under the injury rate is flat across load. Among those aged 29 or over it
runs from 12\% at low load to 48\% at high load.

\section{The model}

Each player's injury is $Y_i \dist \text{Bernoulli}(\pi_i)$ with
\[
  \log\frac{\pi_i}{1 - \pi_i} = \bm{x}_i^\top \bm{\beta},
\]
so a coefficient is a change in the log odds and $e^{\beta}$ is an odds
ratio. Age is centred at 25 and training load at 3, which makes the intercept
and main effects refer to a typical player. \mintinline{r}|minutes_played| is
excluded for the reason given in week 3.

\begin{minted}{r}
fit_all <- glm(injured ~ position + age_c + load_c + prev_injuries + bmi,
               family = binomial, data = players)
out(round(coef(summary(fit_all)), 4), "coef-all")

fit_main <- glm(injured ~ age_c + load_c + prev_injuries,
                family = binomial, data = players)
drop_test <- anova(fit_main, fit_all, test = "Chisq")

add_term <- function(term) {
  bigger <- update(fit_main, as.formula(paste(". ~ . +", term)))
  test <- anova(fit_main, bigger, test = "Chisq")
  data.frame(deviance = test$Deviance[2], df = test$Df[2],
             p = test$`Pr(>Chi)`[2], row.names = paste("add", term))
}
tests <- rbind(
  data.frame(deviance = drop_test$Deviance[2], df = drop_test$Df[2],
             p = drop_test$`Pr(>Chi)`[2], row.names = "drop position, bmi"),
  add_term("age_c:load_c"),
  add_term("age_c:prev_injuries"),
  add_term("load_c:prev_injuries"),
  add_term("team")
)
out(round(tests, 4), "tests")
\end{minted}

\inputminted{text}{output/coef-all.txt}
\inputminted{text}{output/tests.txt}

Each row is a likelihood-ratio test: the change in deviance between nested
models is compared with $\chi^2$ on the change in degrees of freedom.
Position and BMI can be dropped together ($p = 0.55$). Of the three two-way
interactions only age by training load matters, and it matters strongly
(deviance 13.97 on 1 degree of freedom, $p = 0.0002$). Adding the 23 club
indicators does not help ($p = 0.59$).

\begin{minted}{r}
fit_inj <- update(fit_main, . ~ . + age_c:load_c)
out(summary(fit_inj), "summary-inj")
out(round(exp(cbind(odds_ratio = coef(fit_inj), confint(fit_inj))), 3),
    "odds-ratios")
\end{minted}

\inputminted{text}{output/summary-inj.txt}
\inputminted{text}{output/odds-ratios.txt}

The intervals are profile-likelihood intervals. The model lowers the deviance
from 594.7 to 525.3 using four parameters.

\section{Interpretation}

\paragraph{Previous injuries.}
Each previous injury multiplies the odds of injury by 1.69 (95\% interval
1.41 to 2.03), whatever the player's age and load.

\paragraph{Training load and age.}
With an interaction in the model, the main effects describe one slice only:
\mintinline{r}|load_c| is the load effect at age 25 and
\mintinline{r}|age_c| the age effect at load 3. The log odds ratio for one
more point of load at age $a$ is
$\hat\beta_{\text{load}} + (a - 25)\hat\beta_{\text{age:load}}$, with a
standard error that needs both variances and their covariance.

\begin{minted}{r}
b <- coef(fit_inj)
v <- vcov(fit_inj)
load_effect <- function(age) {
  k <- age - 25
  estimate <- b[["load_c"]] + k * b[["age_c:load_c"]]
  se <- sqrt(v["load_c", "load_c"] + k^2 * v["age_c:load_c", "age_c:load_c"] +
               2 * k * v["load_c", "age_c:load_c"])
  data.frame(age, odds_ratio = exp(estimate),
             lower = exp(estimate - 1.96 * se),
             upper = exp(estimate + 1.96 * se))
}
out(round(map_dfr(c(20, 25, 30), load_effect), 2), "load-effect")

grid <- expand_grid(age = 17:36, training_load = 1:5, prev_injuries = 0) |>
  mutate(age_c = age - 25, load_c = training_load - 3)
grid$probability <- predict(fit_inj, grid, type = "response")

fig <- ggplot(grid, aes(x = age, y = probability,
                        colour = factor(training_load))) +
  geom_line() +
  scale_colour_manual(values = pal_load, name = "Training load") +
  labs(x = "Age", y = "Fitted probability of injury")
save_fig(fig, "fitted", height = 2.8)
\end{minted}

\inputminted{text}{output/load-effect.txt}

\begin{figure}[H]
  \centering
  \incfig{fitted}
  \caption{Fitted injury probability for a player with no previous injuries.
    The lines cross near age 23: below it load carries no extra risk, above
    it the lines fan out.}
  \label{fig:fitted}
\end{figure}

At 20 the odds ratio per point of load is 0.71 with an interval that includes
1, so there is no evidence of any load effect for young players. At 25 it is
1.24, and at 30 it is 2.17 (1.54 to 3.05). In probabilities, a 30-year-old
with no injury history has a fitted risk of 7\% at load 2 and 42\% at load 5.

\paragraph{Caution.}
The curves below age 23 fall with load because a single product term must
bend both ways. The data there are consistent with no effect, and week 5
checks whether that part of the fit can be trusted.

\section{A second binary outcome: winning}

\begin{minted}{r}
fit_win <- glm(result == "Win" ~ home_away + opponent_rank,
               family = binomial, data = matches)
out(round(cbind(coef(summary(fit_win))[, c(1, 2, 4)],
                exp(cbind(odds_ratio = coef(fit_win), confint(fit_win)))), 3),
    "win")
\end{minted}

\inputminted[fontsize=\footnotesize]{text}{output/win.txt}

Playing at home multiplies Ashford's odds of winning by 4.9, and each place
lower the opponent finished multiplies them by 1.09. With 46 matches the
intervals are wide (1.4 to 19.9 for home advantage), and the opponent effect
does not reach the 5\% level ($p = 0.07$). Reducing a match to win or not
discards the score, so weeks 6 and 10 model the goals and the three-level
result.

\section{Summary}

\begin{itemize}
  \item Final model: \mintinline{r}|injured ~ age_c * load_c + prev_injuries|
    with a logit link.
  \item Previous injuries act the same for everyone. The evidence that training
    load is a risk factor is confined to older players.
  \item The residual deviance of a binary model cannot test fit, because each
    observation is 0 or 1. Week 5 checks this model with tools made for
    binary data.
\end{itemize}

\end{document}
